\setchapterpreamble[u]{\margintoc}
\chapter{数学}
\labch{mathematics}

数学排版是 \LaTeX{} 的看家本领，也是许多用户选择 \LaTeX{} 的首要原因。
本章将展示 \Class{kaobook} 中文模板中的定理环境、公式排版以及矩阵等数学功能。

\index{数学排版}

\section{定理环境概述}

\Class{kaobook} 通过 \Package{kaotheorems} 宏包提供了一套完整的定理环境体系。
这些环境共享一个浅黄色背景的彩色盒子（通过 \Package{tcolorbox} 实现），
使得数学定义和定理在视觉上从正文中凸显出来\sidenote{定理环境的背景颜色可以通过
\texttt{\textbackslash usepackage[framed,background=red!15!white]\{kaotheorems\}}
来自定义。每种环境还可以单独配颜色！}。

所有定理环境按编号方式分为两类：

\begin{description}
    \item[共享计数器] Theorem、Proposition、Lemma、Corollary —— 这四个环境
        共享同一个计数器，按节（section）重置。例如，如果节 2.1 中有一个 Theorem 2.1.1，
        那么下一个 Proposition 就是 Proposition 2.1.2。
    \item[独立计数器] Definition、Assumption、Remark、Example、Exercise ——
        这些环境各自拥有独立的计数器，同样按节重置。
\end{description}

\index{定理环境}

\section{定义与定理}

让我们从一个经典的定义开始：

\begin{definition}[度量空间]
\labdef{metric-space}
设 $X$ 是一个非空集合。若函数 $d: X \times X \to \mathbb{R}$ 满足以下条件，
则称 $(X, d)$ 为\textbf{度量空间}\index{度量空间}，$d$ 为 $X$ 上的\textbf{度量}：
\begin{enumerate}[label=(\roman*)]
    \item 非负性：$\forall x, y \in X,\; d(x, y) \ge 0$，且 $d(x, y) = 0 \iff x = y$；
    \item 对称性：$\forall x, y \in X,\; d(x, y) = d(y, x)$；
    \item 三角不等式：$\forall x, y, z \in X,\; d(x, z) \le d(x, y) + d(y, z)$。
\end{enumerate}
\end{definition}

\marginnote{度量空间是泛函分析和拓扑学的基础概念。欧几里得空间
$\mathbb{R}^{n}$ 上的标准距离 $d(x, y) = \sqrt{\sum (x_i - y_i)^{2}}$
是最常见的度量空间例子。}

\refdef{metric-space} 是分析学中最重要的基本概念之一\index{分析学}。
基于此定义，我们可以陈述一系列基本定理：

\begin{theorem}[开集的性质]
\labthm{open-set-prop}
设 $(X, d)$ 是一个度量空间。
\begin{enumerate}[label=(\roman*)]
    \item 任意多个开集的并集仍然是开集；
    \item 有限多个开集的交集仍然是开集；
    \item $X$ 和 $\emptyset$ 都是开集。
\end{enumerate}
\end{theorem}

\sidenote[][*-1]{开集的定义请参阅 \refdef{open-set}（稍后定义）。
这是度量空间理论中最基础的定义之一。}

\begin{definition}[开集]
\labdef{open-set}
设 $(X, d)$ 为度量空间。子集 $U \subseteq X$ 称为\textbf{开集}\index{开集}，
如果对于任意 $x \in U$，都存在 $\varepsilon > 0$ 使得开球
$B(x, \varepsilon) = \{y \in X : d(x, y) < \varepsilon\}$ 包含于 $U$ 中。
\end{definition}

\begin{proposition}[柯西序列的有界性]
\labprop{cauchy-bounded}
在度量空间 $(X, d)$ 中，每个柯西序列都是有界的。
\end{proposition}

\begin{proof}
设 $\{x_n\}$ 是 $X$ 中的柯西序列。取 $\varepsilon = 1$，则存在 $N \in \mathbb{N}$，
使得当 $m, n \ge N$ 时，$d(x_m, x_n) < 1$。令
\[
R = \max\{1, d(x_1, x_N), d(x_2, x_N), \ldots, d(x_{N-1}, x_N)\},
\]
则对于所有 $n \in \mathbb{N}$，有 $d(x_n, x_N) \le R$，即序列有界。
\end{proof}

\index{柯西序列}
\index{有界性}

\begin{lemma}
\lablemma{finite-intersection}
若 $U_1, U_2, \ldots, U_n$ 均为开集，则 $\bigcap_{i=1}^{n} U_i$ 也是开集。
\end{lemma}

\begin{corollary}
\labcor{open-ball}
每个开球 $B(x, \varepsilon)$ 都是开集。特别地，$\mathbb{R}^{n}$ 中的每个开区间都是开集。
\end{corollary}

\begin{remark}
注意 \refthm{open-set-prop} 中「有限多个」和「任意多个」之间的区别。
无限多个开集的交集未必是开集——例如，
$\bigcap_{n=1}^{\infty} (-1/n,\; 1/n) = \{0\}$，而单点集在 $\mathbb{R}$ 中不是开集。
这正是度量空间拓扑丰富性的体现。
\end{remark}

\begin{example}[离散度量空间中的开集]
\labexample{discrete-metric}
在离散度量空间\index{离散度量空间}中，对于任意 $x, y \in X$，
\[
d(x, y) = \begin{cases}
    0, & x = y, \\
    1, & x \neq y.
\end{cases}
\]
此时，$X$ 的\textbf{每一个}子集都是开集。因为对于任意 $x \in U \subseteq X$，
取 $\varepsilon = 1$，开球 $B(x, 1) = \{x\} \subseteq U$。
\end{example}

\section{公式排版}

\Class{kaobook} 支持 \LaTeX{} 的所有标准数学环境。本节展示一些常用公式样式。

\subsection{行内公式与行间公式}

行内公式用 \lstinline|$...$| 包裹：如欧拉恒等式 $e^{i\pi} + 1 = 0$，
它被广泛认为是数学中最优美的公式之一\sidenote{欧拉恒等式将五个最基本的数学常数
（$e$、$i$、$\pi$、$1$、$0$）联系在一起，堪称数学之美的极致体现。
参见 \cite{needham1997}。}。行间公式用 \lstinline|\[ ... \]| 或
\lstinline|\begin{equation} ... \end{equation}|：

\begin{equation}
    e^{i\theta} = \cos\theta + i\sin\theta
    \label{eq:euler}
\end{equation}

\index{欧拉公式}

公式 \refeq{euler} 被称为\textbf{欧拉公式}\index{欧拉公式}，它是复分析中的基石。
当 $\theta = \pi$ 时，即得欧拉恒等式。

\subsection{多行对齐公式}

使用 \lstinline|align| 环境可以对齐多个公式：

\begin{align}
    (x + y)^2 &= x^2 + 2xy + y^2 \label{eq:square-sum} \\
    (x - y)^2 &= x^2 - 2xy + y^2 \label{eq:square-diff} \\
    x^2 - y^2 &= (x + y)(x - y) \label{eq:diff-squares}
\end{align}

\index{代数恒等式}

使用 \lstinline|gather| 环境可以居中排列多个公式：

\begin{gather}
    \int_{0}^{\infty} \frac{\sin x}{x}\,dx = \frac{\pi}{2} \label{eq:dirichlet} \\
    \sum_{n=1}^{\infty} \frac{1}{n^2} = \frac{\pi^2}{6} \label{eq:basel}
\end{gather}

\sidenote[][*-1]{\refeq{dirichlet} 是狄利克雷积分，\refeq{basel} 是巴塞尔问题的解。
两者都深刻地体现了 $\pi$ 在分析学中的核心地位。}

\subsection{分段函数}

分段函数可以用 \lstinline|cases| 环境：

\begin{equation}
    f(x) = \begin{cases}
        x^2 \sin(1/x), & x \neq 0, \\
        0,              & x = 0.
    \end{cases}
    \label{eq:piecewise}
\end{equation}

\refeq{piecewise} 是一个经典的反例：它处处可导，但导数在 $x = 0$ 处不连续。

\index{分段函数}

\section{矩阵}

\subsection{基础矩阵}

矩阵\index{矩阵}可以通过多种环境排版。最常用的是 \lstinline|pmatrix|（圆括号矩阵）：

\begin{equation}
    A = \begin{pmatrix}
        a_{11} & a_{12} & \cdots & a_{1n} \\
        a_{21} & a_{22} & \cdots & a_{2n} \\
        \vdots & \vdots & \ddots & \vdots \\
        a_{m1} & a_{m2} & \cdots & a_{mn}
    \end{pmatrix}
    \label{eq:matrix-basic}
\end{equation}

\subsection{不同类型的矩阵分隔符}

\begin{itemize}
    \item \texttt{pmatrix}：圆括号 $\begin{pmatrix} a & b \\ c & d \end{pmatrix}$
    \item \texttt{bmatrix}：方括号 $\begin{bmatrix} a & b \\ c & d \end{bmatrix}$
    \item \texttt{Bmatrix}：花括号 $\begin{Bmatrix} a & b \\ c & d \end{Bmatrix}$
    \item \texttt{vmatrix}：行列式 $\begin{vmatrix} a & b \\ c & d \end{vmatrix}$
    \item \texttt{Vmatrix}：双竖线 $\begin{Vmatrix} a & b \\ c & d \end{Vmatrix}$
\end{itemize}

\marginnote[*-18]{矩阵维度过大时，可以配合 \texttt{\textbackslash scalebox} 或 
\texttt{\textbackslash medmath}（需要 \texttt{nccmath}）进行缩放。}

\subsection{增广矩阵与分块矩阵}

增广矩阵\index{增广矩阵}通常使用 \lstinline|array| 环境：

\begin{equation}
    \left(\begin{array}{ccc|c}
        1 & 2 & 3 & 4 \\
        0 & 1 & 5 & 6 \\
        0 & 0 & 1 & 7
    \end{array}\right)
    \label{eq:augmented}
\end{equation}

\section{综合练习}

让我们通过一个练习题来巩固以上概念：

\begin{exercise}
\labexercise{intro-exercise}
考虑度量空间 $(\mathbb{R}^2, d)$，其中 $d$ 是欧几里得距离。
\begin{enumerate}[label=(\arabic*)]
    \item 证明（或举反例）：单位开球 $B((0,0), 1)$ 不能写成两个开集的不交并，
        除非其中一个是空集。
    \item 计算下列矩阵的行列式：
    \[
    \begin{vmatrix}
        1 & 2 & 3 \\
        2 & 3 & 1 \\
        3 & 1 & 2
    \end{vmatrix}
    \]
    \item 在 $\mathbb{R}$ 上考虑度量 $d(x, y) = |\arctan x - \arctan y|$。
        $(\mathbb{R}, d)$ 是否完备？请说明理由。
\end{enumerate}
\index{欧几里得距离}
\end{exercise}

\sidenote{关于完备度量空间（Complete Metric Space）的更多讨论，
请参考 Rudin 的 \textit{Principles of Mathematical Analysis}\cite{rudin1976}。}

\begin{kaobox}[title=定理环境速查表]
下表汇总了 \Class{kaobook} 提供的所有数学定理环境：
\begin{center}
\begin{tabular}{l l l}
    \toprule
    环境名 & 名称 & 编号方式 \\
    \midrule
    \texttt{definition}  & Definition  & 独立计数器 \\
    \texttt{theorem}     & Theorem     & 共享计数器 \\
    \texttt{proposition} & Proposition & 共享计数器 \\
    \texttt{lemma}       & Lemma       & 共享计数器 \\
    \texttt{corollary}   & Corollary   & 共享计数器 \\
    \texttt{assumption}  & Assumption  & 独立计数器 \\
    \texttt{remark}      & Remark      & 独立计数器 \\
    \texttt{example}     & Example     & 独立计数器 \\
    \texttt{exercise}    & Exercise    & 独立计数器 \\
    \bottomrule
\end{tabular}
\end{center}
\end{kaobox}

\index{定理环境速查}

数学是语言的延伸，而排版是思想的容器。有了 \Class{kaobook} 强大的定理环境
和公式排版功能，你可以将精力集中在数学内容本身，而不必为排版格式分心。

\endinput
